\Subhead{1.4.1}{Метод гармонического баланса}
Аналитические методы находят приближённое решение; качественный анализ выясняет, какие движения возможны и как они меняются. \textbf{Гармонический баланс} ищет периодическое движение в виде конечной суммы гармоник. Её подставляют в уравнение и приравнивают коэффициенты выбранных синусов и косинусов. Нелинейность создаёт и отброшенные гармоники, поэтому нужно объяснить, почему ими можно пренебречь.

Для слабонелинейного осциллятора достаточно $x=A\cos\omega t$. В свободной модели Дуффинга $\ddot x+\omega_0^2x+\alpha x^3=0$ используем $\cos^3\theta=(3\cos\theta+\cos3\theta)/4$. Баланс первой гармоники даёт
\begin{equation}
 -\omega^2A+\omega_0^2A+\tfrac34\alpha A^3=0,
 \qquad \omega\simeq\omega_0+\frac{3\alpha A^2}{8\omega_0}.
\end{equation}
Цель --- зависимость частоты от амплитуды; для вынужденного режима аналогичный расчёт даёт \eqref{eq:duffingresponse}. При $|\alpha|A^2/\omega_0^2\ll1$ третья гармоника мала и нерезонансна. В добротном генераторе контур дополнительно выделяет основную гармонику. Баланс определяет кандидатов на периодические решения, но сам не определяет их устойчивость или процесс установления. Для релаксационных импульсов одной гармоники недостаточно.
